\section{Two-photon exchange}\label{sec:TPE}

The TPE contribution to lepton scattering, shown in \Figref{TPE_diagram}, is of particular importance for two main reasons: Firstly,  hadronic corrections, cf.\ the blob in \Figref{TPE_diagram} (bottom panel), are notoriously difficult to calculate and often have a large relative uncertainty. Secondly, as will be explained in Sec.~\ref{sec:RoPTcomp}, the OPE may not be a good approximation for the extraction of FFs from the unpolarized cross section at large $Q^2$. Therefore, the leading TPE correction 
\begin{equation}
\label{delta2gTPE}
\delta_{2\gamma}   \simeq   \frac{2\re \big(\AMP^{(0)*} \AMP_{2\gamma}\big)} {|\AMP^{(0)}|^2},
\end{equation}
following from the interference of the Born and TPE amplitudes, $\AMP^{(0)}$ and $\AMP_{2\gamma}$, has to be taken into account appropriately. A complete calculation should go beyond the soft-photon approximation and include the hard TPE with all possible intermediate states. In the last decade, predictions of the TPE correction have been systematically improved and model dependence has been reduced by employing dispersion relations and effective field theories. 


This part of the paper is organized as follows. In Sec.~\ref{sec:RoPTcomp}, we illustrate the importance of the TPE by comparing Rosenbluth and polarization-transfer measurements of the proton FFs. In Sec.~\ref{sec:TPEtheory}, we review theoretical predictions for the leading TPE correction, distinguishing between proton and inelastic intermediate states, as well as regions from small to large momentum transfer. In Sec.~\ref{sec:EFTCalc}, we discuss effective field theory calculations. In Sec.~\ref{sec:TPEempirical}, empirical extractions of TPE corrections and amplitudes are presented and updated based on new data.  In Sec.~\ref{Sec:Experiments}, results of past experiments  aimed at extracting the TPE  are presented and compared to theoretical predictions. We finish with an outlook on future experiments and theory advances in Sec.~\ref{sec:TPEOutlook}. For further reading, we refer to the following reviews in Refs.~\cite{Carlson:2007sp, Arrington:2011dn,Afanasev:2017gsk,Borisyuk:2019gym} and the recent CFNS whitepaper, Ref.~\cite{Afanasev:2020hwg}.

\subsection{Rosenbluth vs. polarization transfer experiments}\seclab{RoPTcomp}




Polarization experiments in elastic $ep$ scattering at Jefferson Lab Hall A \cite{JeffersonLabHallA:1999epl,Gayou:2001qt,Punjabi:2005wq,Puckett:2011xg}, revealed that the ratios of proton FFs, $G_E(Q^2)/G_M(Q^2)$, extracted based on the Rosenbluth \cite{Rosenbluth:1950yq} or polarization transfer methods \cite{Akhiezer:1968ek,Akhiezer:1974em} in the OPE approximation deviate with increasing $Q^2$. The FF ratio from Rosenbluth extractions is nearly constant as a function of $Q^2$, while it decreases for polarization transfer measurements, as can be seen in Fig.\ \ref{fig.GEGM}. In the limit of OPE, one defines an unpolarized reduced cross section that is linear in $\varepsilon$ and the corresponding ratio of measurements with transverse or longitudinal polarization of the recoiling proton is constant in $\varepsilon$  as given in Eqs.~\ref{eq:reducedCS} and \ref{eq:PTOPE}. What is measured in experiment, however, will always include radiative corrections. Taking these into account and under some assumptions, in presence of TPE, the expressions modify according to Eqs.~\ref{eq:delta_TPE_massive} and \ref{PtPlradcorr}, respectively.
One can see that after inclusion of the TPE correction, the interpretation of what is measured in Rosenbluth and polarization transfer experiments changes.  
Model independent considerations  \cite{Rekalo:2003xa,Rekalo:2003km,Rekalo:2004wa} show that it is still possible to recover experimentally the electric and magnetic proton FFs, even in presence of TPE, but this requires either the measurement of three time-odd or five time-even polarization observables (including triple spin observables, of the order of $\alpha$) or, alternatively, the generalization of the Akhiezer-Rekalo recoil proton polarization method  \cite{Akhiezer:1968ek,Akhiezer:1974em} with longitudinally polarized electrons and positron beams in identical kinematical conditions.


In a series of papers in the early 2000's, it has been shown that the inclusion of ``hard'' TPE may reconcile the FF extractions from polarized and unpolarized $\ell p$ scattering, within some assumptions (real-valued FFs and  $\varepsilon$ linearity of the TPE contribution). Guichon and Vanderhaeghen identified in Ref.~\cite{Guichon:2003qm} the experimental observables used to extract the $G_E/G_M$ ratio as the $\varepsilon$-slope in the reduced cross section
\begin{align}
\sigma_R(\varepsilon, Q^2)
= G_{M}^2 \left(1 + \frac{\varepsilon}{\tau}  
R_{EM}^2 \right) \left\{ 1 + \delta_{2 \gamma}\right\} .
\label{eq:sigmaR}
\end{align}
Neglecting lepton mass terms and using the shorthand notations for the TPE amplitudes of 
Eq.~\ref{eq:2gaampl}, $\delta_{2\gamma}$ 
has the form
\begin{align}
&\delta_{2 \gamma}(\varepsilon, Q^2) = \left(1 + \frac{\varepsilon}{\tau}  R_{EM}^2 \right)^{-1} 
\nn \\ &
\times \left\{ 2 \, Y_M
+ 2 \varepsilon  \frac{R_{EM}}{\tau} \, Y_E + 2 \varepsilon \left( 1 + \frac{R_{EM}}{\tau} \right) Y_3 \right\}.
\label{crossen} 
\end{align} 
On the other hand, the polarization transfer experiments yield the ratio given by Eq.~\ref{PtPlradcorr}
\begin{align}
-\sqrt{\frac{\tau (1+\varepsilon)}{2\varepsilon}}  \frac{P_t}{P_l}&= R_{EM} +Y_E - R_{EM} Y_M \nn\\&  +\Bigg (1-\frac{2 \varepsilon}{1+\varepsilon} R_{EM} \Bigg )Y_3.
\eqlab{Rratios}
\end{align}
Reference \cite{Guichon:2003qm} showed that the TPE in the Rosenbluth method effectively corrects a small number, cf.\ the 
$\varepsilon$ dependent term in Eq.~\ref{eq:sigmaR} with $R_{EM}^2(0)\sim (1/2.79)^2 \sim 0.13$ which has an additional $1/\tau$ suppression at large $Q^2$. The FF ratio extracted from $P_t/P_l$ is only mildly affected by radiative corrections as it involves asymmetries,  whereas the cross sections themselves are strongly modified by radiative corrections at moderately and large momenta.  
Around the same time, Blunden et al.\  in Ref.~\cite{Blunden:2003sp} showed with a simple hadronic model calculation, including the finite size of the proton, that the dominant TPE with a proton intermediate state allows to partially resolve the discrepancy.
Inclusion of the $\Delta(1232)$ resonance further improved the agreement, as shown by Kondratyuk et al.\ in Ref.~\cite{Kondratyuk:2005kk}. A similar conclusion was drawn by Chen et al.\ and Afanasev et al.\ in Refs.~\cite{Chen:2004tw,Afanasev:2005prd} respectively, with a partonic model calculation. 


\begin{figure}[ht]
\includegraphics[width=\columnwidth]{TPE/fig/GEGMratio.pdf}
\caption{Ratio of proton electric and magnetic FFs, $G_E \,\mu_p/G_M$, with $\mu_p$ the magnetic moment of the proton, as obtained from two types of measurements: Rosenbluth separation \cite{LITT197040,BARTEL1973429,Andivahis:1994rq,PhysRevD.49.5671,PhysRevC.70.015206,Qattan:2004ht,Christy:2021snt} and polarization transfer \cite{Gayou:2001qt,Punjabi:2005wq,ResonanceSpinStructure:2006oim,Paolone:2010qc,Puckett:2011xg,Puckett:2017flj,SANE:2018cub}. Shown are a subset of the available data to not overcrowd the figure. The Rosenbluth result only contains the TPE corrections in the form of the well-known Maximon-Tjon radiative corrections \cite{Maximon:2000hm} (no hard TPE was included), whereas for the polarized extraction no radiative corrections were taken into account. }
\label{fig.GEGM}
\end{figure}

 Note that alternative explanations have been put forward, such as different calculations of radiative corrections, including also  lepton structure functions \cite{Gramolin:2016hjt,Bystritskiy:2006ju}, correlations between the Rosenbluth parameters \cite{Tomasi-Gustafsson:2006dhf},
or acceptance problems in the analysis or experiment \cite{Pacetti:2016tqi}.


The electric FF is determined by the slope of the Rosenbluth plot, at fixed $Q^2$, derived from the radiatively corrected cross section. Corrections at large $Q^2$ may reach 40\% \cite {Andivahis:1994rq} and the $\varepsilon$-slope may essentially change. Even more dramatically, above 3 GeV$^2$ the slope of the measured cross section may even be negative before corrections, cf.\ Fig.~\ref{Fig:AndiRad} \cite{Tomasi-Gustafsson:2004fss,Perdrisat:2006hj} where the effect of radiative corrections \cite{Mo:1968cg}, including ``soft'' but no ``hard'' TPE, is shown. Note that $G_E$ must be real in the space-like region. This means that $G_E^2$ is essentially determined by the $\varepsilon$-dependence of the applied radiative corrections. Different calculations show a difference of few percent in the slope, already at NLO \cite{Gramolin:2014pva,Gerasimov:2015aoa}.



\begin{figure}[ht]
\includegraphics[width=\columnwidth]{TPE/fig/Fig10.eps}
\caption{Rosenbluth plot for the data from Ref.~\cite{Andivahis:1994rq} at momentum transfers of 1.75 (green squares), 3.25 (blue triangles) and 5 GeV$^2$ (red circles). Empty and full symbols show data before and after radiative corrections \cite{Mo:1968cg} (no hard TPE included), respectively. Figure taken from Ref.~\cite{Perdrisat:2006hj}.
}
\label{Fig:AndiRad}
\end{figure}
Concerning polarization observables, it is commonly assumed that the FF ratio given by polarization measurements is not (or less) affected by radiative corrections, giving therefore a more reliable result on this ratio than the one extracted by the Rosenbluth method. 
The polarization asymmetries $P_t$ and $P_l$ are ratios of cross sections, cf.\ Eqs.~\ref{eq:Pl} and \ref{eq:Pt}, in which the bulk of the radiative corrections that factorize (virtual corrections on lepton side and soft-photon emission corrections) drop out. Therefore, they are less sensitive to radiative corrections than the unpolarized cross section.




Already in the 1970's, a reason why the TPE correction could become important at large $Q^2$ was put forward \cite{Gunion:1972bj,Boitsov:1972if,Franco:1973uq}:
when the transferred momentum is equally shared between the two photons the scaling in $\alpha$ may be compensated by the steep decrease of the FFs with $Q^2$. The effect is therefore expected to increase with $Q^2$ and with the hadron mass.
The advent of the high duty cycle and highly polarized electron beam at the Jefferson Lab, together with the availability of large solid angle spectrometers and hadron polarimetry in the GeV region, opened the way to high-precision experiments in elastic and inelastic electron-hadron scattering. Two experiments of elastic electron-deuteron ($ed$) scattering at large $Q^2$ in Hall A \cite{JeffersonLabHallA:1998xrv} and Hall C \cite{JeffersonLabt20:1998vrh} claiming an error of 5\%, showed a discrepancy of up to 15\% in the elastic $ed$ cross section at the same $Q^2$, but different energies and angles. A possible explanation was brought up that the TPE contribution could be at the origin of these findings \cite{Rekalo:1999mt}. Finally the discrepancy was attributed to a systematic error of the Hall C spectrometer setting, as no $Q^2$ dependence was observed.


\subsection{Theoretical predictions} \seclab{TPEtheory}


\begin{figure}[ht]
\centering
\includegraphics[scale=0.45,angle=0]{TPE/fig/ElasticTPE.pdf}\\
\vspace{0.5cm}
\includegraphics[scale=0.45,angle=0]{TPE/fig/InelasticTPE.pdf}
\caption{Off-forward TPE contribution to elastic $\ell p$ scattering with proton intermediate state (top panel), where the circles represent the proton FFs, and inelastic intermediate states (bottom panel), where the blob represents all possible excitations. The horizontal lines correspond to the lepton and proton (double line). Crossed TPE  diagrams are implied in all references to this figure. \label{fig:TPE_diagram}}
\label{OffTPE}
\end{figure}

In this subsection, we review theoretical predictions for the TPE contributions to $\ell p$ scattering. We will refer to the (standard) TPE approximations by Maximon-Tjon \cite{Maximon:2000hm} (or Mo-Tsai \cite{Mo:1968cg}), conventionally included in the experimental analysis, as ``soft'' TPE. (Non-standard) refinements of the TPE, that need to be included given the accuracy of modern scattering experiments, will be referred to below as ``hard'' TPE. Starting from the so-called elastic TPE with a proton intermediate state, Fig.~\ref{OffTPE} (top panel), we subsequently discuss TPE contributions with inelastic intermediate states, Fig.~\ref{OffTPE} (bottom panel), in regions of small, medium and large momentum transfer. Furthermore, we discuss the forward TPE as the limiting factor in the theoretical description of the energy spectra of light muonic atoms.

\def\Re{{\rm Re}}
\subsubsection{Infrared subtraction schemes}\seclab{IRsubtraction}

When plotting the elastic TPE contribution to $\delta_{2\gamma}$ or $R_{2\gamma}$, defined in Eqs.~\ref{delta2gTPE} and \ref{eq:R2gammadef}, one has to somehow treat the IR divergences they contain.
Of course the only way that leads to a measurable quantity is to combine them with the corresponding real corrections as discussed in Sec.~\ref{Sec:RadCorrAdrian}.
However, this can be quite cumbersome in practise and makes comparing different TPE scenarios and parametrisations needlessly complicated.
Instead, one defines an IR subtraction scheme that removes the IR divergence (and potentially some finite parts as well).
The exact form of this subtraction term is somewhat arbitrary and solely determined by conventions since the result is no longer a physical quantity.

A common way of doing this in the high-energy community is the definition of the Catani $I$ operator~\cite{Catani:1998bh} that simply removes the $1/\epsilon$ poles in dimensional regularisation.

In the QED community it is not uncommon to add the first terms of Eqs.~\ref{slicing} or \ref{subtraction}
\begin{eqnarray}
    \delta_{2\gamma}^{\rm eik} &= \delta_{2\gamma} +&
          \int_{E_\gamma < \delta} \dd\Phi_\gamma\ \mathcal{E} \nn\\
    & = \delta_{2\gamma} +&
          \underbrace{
              \int \dd\Phi_\gamma\Big(\frac{p_1\cdot k_1}{p_1\cdot k\ k_1\cdot k}+\frac{p_2\cdot k_2}{p_2\cdot k\ k_2\cdot k}}_{\hat{\mathcal{E}}(s)}\nn\\
              &&
         \hspace*{1cm} - (s\rightarrow u)\Big).
\end{eqnarray}
This scheme, sometimes called eikonal subtraction~\cite{Engel:2019nfw}, has the added advantage of giving the remnant a physical interpretation as long as the parameter $\delta$ is chosen small enough -- it just approximately includes real corrections up to the cut-off $\delta$.
The calculation of this integral is somewhat involved but the result is well known~\cite{Engel:2019nfw,Frederix:2009yq,Ulrich:2020phd}.
Since this physical interpretation exists, it also means that the result needs to be independent of how the TPE calculation was regularised, making comparisons between different methods easier.
The TPE community usually follows either Mo-Tsai (MoT) \cite{Mo:1968cg} or Maximon-Tjon (MTj) \cite{Maximon:2000hm}.
Note that in this paper, we use the MTj convention for our Figs.~\ref{fig.Dsig}-\ref{fig.NearForw}. The latter defines the IR-subtracted TPE contribution as
\begin{align}    &\delta_{2\gamma}^{\rm MTj} = \delta_{2\gamma} -
        \frac{2 \alpha}{\pi} q^2 \,\Re\Bigg[p_1\cdot k_1 \nn\\
        &\times\int  \frac{\dd \ell}{
            \ell^2
            (\ell-q)^2
            \Big[(\ell-k_1)^2-m^2\Big]
            \Big[(\ell+p_1)^2-M^2\Big]
        } \nn\\&\qquad\quad - (s\rightarrow u)\Bigg].
\end{align}
This integral can be evaluated both using photon-mass regularisation or dimensional regularisation
    \begin{eqnarray}
    \delta_{2\gamma}^{\rm MTj} &=& \delta_{2\gamma} +
    \frac{2 \alpha}{\pi}\,\Re\Bigg[\nn\\
        \frac{m^2+M^2-s}{\sqrt{\lambda_s}}
        & \times&\ln\frac{m^2+M^2-s + \sqrt{\lambda_s}}{2mM}
    \begin{cases} \ln\frac{m_\gamma^2}{Q^2} \\ \frac1{\epsilon}\Big(\frac{\mu^2}{Q^2}\Big)^\epsilon \end{cases}\nn\\
    &
    &- (s\rightarrow u)\Bigg],
\end{eqnarray}
where we have defined $\lambda_s = m^4 + (M^2 - s)^2 - 2 m^2 (M^2 + s)$.
In the limit of $m^2\ll s,M^2,Q^2$, the expression for $e^- p$ scattering simplifies to
    \begin{equation}
    \delta_{2\gamma}^{\rm MTj} = \delta_{2\gamma} +
    \frac{2 \alpha}{\pi}
        \ln\Big|\frac{s - M^2}{u - M^2}\Big|
    \begin{cases} \ln\frac{m_\gamma^2}{Q^2} \\ \frac1{\epsilon}\Big(\frac{\mu^2}{Q^2}\Big)^\epsilon \end{cases}\!\!\!\!\!\!
    .
\end{equation}
To derive the MoT subtraction prescription, we instead start by replacing either the $\ell^2$ propagator or the $(\ell-q)^2$ propagator with $1/q^2$.
The idea is to mimic the soft behaviour of the integral as either the $\ell^2$ propagator goes to zero or the $(\ell-q)^2$ propagator.
This procedure results in two identical triangle functions
    \begin{align}
    &\delta_{2\gamma}^{\rm MoT'} = \delta_{2\gamma} -
        \frac{2\alpha}{\pi} q^2 \Re\Bigg[p_1\cdot k_1 \nn\\
      &  \hspace{-0.5cm}\times\int  \frac{\dd \ell}{
            \ell^2
            q^2
            \Big[(\ell-k_1)^2-m^2\Big]
            \Big[(\ell+p_1)^2-M^2\Big]
        } \nn\\&\qquad\quad - (s\rightarrow u)\Bigg].
        \label{eq:motorgi}
\end{align}
Note that this is not yet the correct MoT prescription, hence {\rm MoT$'$}.
The loop integral is slightly more involved than the previous one but can still be trivially written in terms of the well-known Ellis-Zanderighi functions~\cite{Ellis:2007qk}.
In this case, we need the integral $I_6^{\rm fin}(s;m^2,M^2)$ commonly referred to as Triangle 6,
\begin{eqnarray}
    \delta_{2\gamma}^{\rm MoT'} &=& \delta_{2\gamma} +
        \frac{2 \alpha}{\pi} \Bigg\{
            (m^2 + M^2 - s) \Re\Bigg[
                I_6^{\rm fin}(s;m^2,M^2) \nn\\
                &+& \frac1{\sqrt{\lambda_s}}\ln\frac{m^2+M^2-s + \sqrt{\lambda_s}}{2mM} 
            \begin{cases} \ln\frac{m_\gamma^2}{mM} \\ \frac1{\epsilon}\Big(\frac{\mu^2}{mM}\Big)^\epsilon \end{cases}\!\!\!\!\!\!\Bigg]\nn\\
        &-& (s\rightarrow u)\Bigg\}.
\end{eqnarray}
If one chooses, $I_6^{\rm fin}$ could be expressed in terms of logarithms and dilogarithms.
The resulting expression is not overly complicated but reproducing it here serves no practical purpose.
However, in the limit of $m^2\ll s,M^2,Q^2$ the expression is fairly compact
\begin{eqnarray}
    I_6^{\rm fin}(s;m^2,M^2) &=& -\frac{y}{2M^2}\Bigg(
        -\frac13\big(i\pi+\log(-y^3)\big)^2\nn\\ &+ &2\ln(-y) \ln\frac y{1+y} - 2{\rm Li}_2(-y) \nn\\&-& \ln(-y)\ln\frac{m^2}{M^2}
    \Bigg) + \mathcal{O}(m^2)\,,
\end{eqnarray}
where we have defined $y=M^2/(s-M^2)$ and used the usual definitions of the $\zeta$ function and dilogarithm.

To arrive at the correct MoT prescription~\cite{Mo:1968cg} (cf. also~\cite{Arrington:2011dn}) we need to flip $k_1\to-k_1$ in the first term of~\eqref{eq:motorgi}, resulting in
\begin{align}
    &\delta_{2\gamma}^{\rm MoT} = \delta_{2\gamma} -
        \frac{2\alpha}{\pi} q^2 \Re\Bigg[\nn\\
      &  \hspace{-0.5cm}\phantom{+}\int \dd \ell \frac{-p_1\cdot k_1 }{
            \ell^2
            q^2
            \Big[(\ell+k_1)^2-m^2\Big]
            \Big[(\ell+p_1)^2-M^2\Big]
        } \nn\\
      &  \hspace{-0.5cm}+\int \dd \ell \frac{-p_1\cdot k_2 }{
            \ell^2
            q^2
            \Big[(\ell+k_2)^2-m^2\Big]
            \Big[(\ell+p_1)^2-M^2\Big]
        } \Bigg].
\end{align}
This operation is equivalent to setting $s\to s'=2 m^2 + 2 M^2 - s$, which spoils the crossing symmetry but eliminates the $(i\pi)^2$ from the analytic continuation of the logarithm.
The resulting expression can still be written in terms of $I_6^{\rm fin}$ but now in terms of $s'$ rather than $s$
\begin{align}
    \delta_{2\gamma}^{\rm MoT} &= \delta_{2\gamma} +
        \frac{2 \alpha}{\pi} \mathrm{Re}\Bigg [
              (s - m^2 - M^2) I_6^{\rm fin}(s';m^2,M^2) \nn\\
            &\qquad+ (u - m^2 - M^2) I_6^{\rm fin}(u;m^2,M^2)\nn\\
            &+\Bigg[
                  \frac{s-m^2-M^2}{\sqrt{\lambda_{s}}}\ln\frac{-m^2-M^2+s + \sqrt{\lambda_{s}}}{2mM}\nn\\
               &+\frac{u-m^2-M^2}{\sqrt{\lambda_{u}}}\ln\frac{ m^2+M^2-u + \sqrt{\lambda_{u}}}{2mM}
            \Bigg]\nn\\
            &\begin{cases} \ln\frac{m_\gamma^2}{mM} \\ \frac1{\epsilon}\Big(\frac{\mu^2}{mM}\Big)^\epsilon \end{cases}\!\!\!\!\!\!\Bigg\} \Bigg].
\end{align}
The symmetry of the $\lambda$ functions means that $\lambda_{s'}=\lambda_s$, which helps to simplify the result.
The difference between $\delta_{2\gamma}^{\rm MoT}$ and $\delta_{2\gamma}^{\rm MoT'}$ is
\begin{eqnarray}
    &\delta_{2\gamma}^{\rm MoT'}-\delta_{2\gamma}^{\rm MoT} = \dfrac{2\alpha}\pi (m^2+M^2-s)\Re\Bigg[\nn\\
            &+I_6^{\rm fin}(s';m^2,M^2) + I_6^{\rm fin}(s;m^2,M^2)\Bigg].
\end{eqnarray}
Crucially, the poles agree after taking the real part so that the IR cancellation in $\delta_{2\gamma}$ is unaffected.

\subsubsection{Proton intermediate state} \label{subsec:proton_state}

The first terms in the low-$Q$ expansion of the TPE with proton intermediate state can be described model independently. The leading term is given by the well-known Feshbach correction
 \cite{McKinley:1948zz}
\begin{equation} \eqlab{Feshbach}
    \delta_{\text{F}}=Z \alpha \pi \frac{\sin \frac{\theta}{2}-\sin^2 \frac{\theta}{2}}{\cos^2 \frac{\theta}{2}},
\end{equation}
which stems from the interaction between a massless lepton and a structureless proton.  Inelastic intermediate states start to contribute in the subleading $Q^2 \ln Q^2$ term only \cite{Brown:1970te}, see discussion in Sec.~\ref{SmallMT}. It follows that TPE could become larger in lepton scattering on heavy ions, inducing  a large charge asymmetry even at small angles \cite{Kuraev:2009hj}.

During the last two decades,  predictions of the TPE contribution improved considerably. The simplest approach to evaluate the elastic TPE is with a hadronic model calculation of the box (and crossed-box) diagrams in Fig.~\ref{OffTPE} (top panel), assuming on-shell proton FFs, see for instance Refs.~\cite{Blunden:2003sp} and \cite{Tomalak:2014dja} for electron and muon scattering, respectively. A way to avoid model dependence is to instead use a dispersive approach to describe the scalar amplitudes $\cG_i$ and $\cF_i$, introduced in Eqs.~\ref{eq:str_ampl} and \ref{eq:cGamps}. 
Dispersion relations (DRs) allow one to express the real part of the amplitudes via their imaginary part, and the latter, using unitarity, can be related to physical observables. The $s-$channel cut in the elastic TPE diagram then permits the use of on-shell proton FFs. 
First dispersive evaluations of the elastic TPE in $ep$ scattering, neglecting the electron mass, used unsubtraced DRs for the scalar amplitudes \cite{Borisyuk:2006fh,Borisyuk:2008es}. Considering the case of $\mu p$ scattering, in which the lepton mass cannot be neglected, $\cF_4$ will require a once-subtracted DR \cite{Tomalak:2018jak}.

The contribution from TPE with inelastic intermediate states, discussed in the following, is becoming important with increasing momentum transfer. It has been  suggested that further subtractions, e.g., in the DR of the $\cF_3$ amplitude, can be introduced in order to minimize model-dependence from higher intermediate states \cite{Tomalak:2014sva}. Sum rules, which are exact in QED and approximated in QCD, give indications on the contribution of intermediate states \cite{Kuraev:2006ym}. 

\subsubsection{Zero momentum transfer}\seclab{VVCS}
In the limit of zero momentum transfer, $Q^2=0$, the TPE is not important as a radiative correction to the scattering process. $\delta_{2\gamma}$, as defined by the interference of OPE and TPE amplitudes in Eq.~\ref{delta2gTPE}, is vanishing in the forward limit ($Q^2\rightarrow 0$) at fixed $\nu$, see discussion in Ref.~\cite{Tomalak:2018jak}. At the next order in $\alpha$, the contribution from $\vert \AMP_{2\gamma}\vert^2$ to the cross section at zero momentum transfer is not vanishing, however, numerically suppressed.
On the contrary, the TPE in forward kinematics is important in the description of the spectra of light muonic atoms, where its effect is enhanced as compared to ordinary atoms  due to the heavier muon mass. Let us focus on (muonic) hydrogen. The forward TPE corresponds to a $\delta(\vec{r}\,)$ potential that gives an $\mathcal{O}(\alpha^5)$ contribution to the energy levels. For comparison, the leading contribution from the Coulomb potential is of $\mathcal{O}(\alpha^2)$ and the proton finite-size starts to contribute at $\mathcal{O}(\alpha^4)$ through an OPE diagram with proton FFs. The off-forward TPE starts to contribute at $\mathcal{O}(\alpha^6\ln \alpha)$ through the so-called Coulomb distortion effect. See Ref.~\cite{Pachucki:2022tgl} for a comprehensive theory of the Lamb shift in light muonic atoms. Here, we want to focus on the dispersive formalism used to evaluate the $\mathcal{O}(\alpha^5)$ effect of the forward TPE, since a similar approach is used to approximate the small momentum transfer corrections to the scattering process, cf.\ Sec.~\ref{SmallMT}.
 
The forward TPE between a lepton and a proton can be expressed in terms of forward doubly-virtual Compton scattering (VVCS) off the proton. The TPE-induced shift of the $nS$-level in a hydrogen-like atom is given by the unpolarized VVCS amplitudes $T_1$ and $T_2$  \cite{Carlson:2011zd}
\begin{eqnarray}
&&\Delta E^{2\gamma}(nS)= 8\pi \alpha m \,\phi_n^2\,
\frac{1}{i}\int_{-\infty}^\infty \!\frac{\mathrm{d}\bar \nu}{2\pi} \int \!\!\frac{\mathrm{d} \boldsymbol{q}}{(2\pi)^3} \qquad\\
&&\times\frac{\left(\bar Q^2-2\bar \nu^2\right)T_1(\bar \nu,\bar Q^2)-(\bar Q^2+\bar \nu^2)\,T_2(\bar \nu,\bar Q^2)}{\bar Q^4(\bar Q^4-4m^2\bar \nu^2)},\nonumber
\end{eqnarray}
where $\phi_n^2=1/(\pi n^3 a^3)$ is the wavefunction at the origin, $a=(Z \alpha m_r)^{-1}$ is the Bohr radius (in the following $Z=1$ for hydrogen), $m_r$ is the reduced mass of the lepton-proton system and $m$ is the lepton mass ($m_e$ or $m_\mu$, respectively, for hydrogen and muonic hydrogen). Furthermore, $\bar \nu=q_0$ and $\bar Q^2 = \boldsymbol{q}^2 -q_0^2$ are the energy and virtuality of the VVCS photons inside the TPE loop diagram. Similarly, the TPE contribution to the hyperfine splitting can be expressed through the spin-dependent VVCS amplitudes $S_1$ and $S_2$. The TPE effect can then be evaluated with a data-driven dispersive approach \cite{Pachucki:1999zza,Martynenko:2005rc,Carlson:2011zd,Birse:2012eb,Gorchtein:2013yga,Hill:2016bjv,Tomalak:2018uhr,Faustov:2006ve,Carlson:2008ke,Carlson:2011af}, or predicted from variants of chiral perturbation theory \cite{Alarcon:2013cba,Hagelstein:2015lph,Lensky:2017bwi,Hagelstein:2018bdi,Peset:2014yha,Peset:2014jxa,Peset:2016wjq}, and in the future, from lattice QCD \cite{Fu:2022fgh,Hagelstein:2020awq}. For recent reviews, discussing low-energy proton structure in the context of muonic-hydrogen spectroscopy and lepton scattering, including TPE and the proton radius, we refer to Refs.~\cite{Carlson:2007sp,Carlson:2015jba,Gao:2021sml,Peset:2021iul,Antognini:2022xoo,Hagelstein:2015egb}.

Data-driven evaluations of the forward TPE make use of dispersion relations for the VVCS amplitudes, expressing them through proton structure functions measured in electron scattering
\begin{align}
\eqlab{T1DR}
T_1 (\bar  \nu, \bar Q^2) &=T_1(0,\bar Q^2)  \nn\\
&+\frac{32\pi \alpha M\bar \nu^2}{\bar Q^4}\int_{0}^1 
\,\dd x \, 
\frac{x F_1 (x, \bar Q^2)}{1 - x^2 (\bar \nu/\bar \nu_{\mathrm{el}})^2},\\
\eqlab{T2DR}
T_2 (\bar  \nu, \bar Q^2)
&=\frac{16\pi \alpha M}{\bar Q^2} \int_{0}^1 
\!\dd x\, 
\frac{F_2 (x, \bar Q^2)}{1 - x^2 (\bar \nu/\bar \nu_{\mathrm{el}})^2 },
\end{align}
where $F_1(x,\bar Q^2)$ and $F_2(x,\bar Q^2)$ are the unpolarized structure functions, with $x=\bar Q^2/2M\bar \nu$ the Bjorken variable, and $\bar \nu_\mathrm{el}=\bar Q^2/2M$. Similarly, the spin-dependent VVCS amplitudes $S_1$ and $S_2$ follow from dispersion integrals over the spin structure functions $g_1(x,\bar Q^2)$ and $g_2(x,\bar Q^2)$.
Unfortunately, the $T_1$ amplitude in \Eqref{T1DR} requires a once-subtracted dispersion relation. Since the subtraction
function $T_1(0,\bar Q^2)$ cannot be fully constrained from experiment, the data-driven approach suffers from a model dependence. 


The TPE contributions to the energy levels in muonic atoms are usually split into a 
so-called ``elastic'' contribution, corresponding to the TPE with proton intermediate state shown in \Figref{TPE_diagram} (top panel), and a so-called ``polarizability'' contribution, corresponding to the diagrams in \Figref{TPE_diagram} (bottom panel). The elastic part of the VVCS amplitudes is well-known and can be expressed through proton FFs. The uncertainty of the elastic TPE can therefore be reduced with improved input for the proton FFs. This has recently been done based on the FF descriptions from Refs.~\cite{Borah:2020gte,Lin:2021xrc}, which are both consistent with the small proton charge radius extracted from the muonic-hydrogen Lamb shift, see Refs.~\cite{Antognini:2022xqf} and \cite{Pachucki:2022tgl} for the elastic TPE in the muonic hydrogen Lamb shift and hyperfine splitting, respectively. Therefore, the dispersive approach is only used to reconstruct the inelastic part of the VVCS amplitudes.  To this end, the dispersion integral is restricted to the inelastic structure functions, i.e., the region of $x \in [0,x_0]$ with the inelastic threshold $x_0$, e.g., the pion production threshold. In the following section, we will discuss the extension of this formalism to the TPE correction to $\ell p$ scattering in the region of small momentum transfer.


\subsubsection{Small momentum transfer}\label{SmallMT}

In the near-forward kinematics, i.e., in the small momentum transfer region, the TPE amplitude can be approximated with a modification of the dispersive approach presented above, see Eqs.~\ref{eq:T1DR} and \ref{eq:T2DR}, that reconstructs the VVCS amplitudes based on empirical input for the inelastic proton structure functions. In the following, we review the leading-$Q^2$ behaviour of the inelastic TPE corrections to the unpolarized cross section and the beam normal spin asymmetry as functions of the transverse cross section for photoabsorption off the proton. Recall that the leading terms in the $Q^2$ expansion of the TPE stem from the proton intermediate state \cite{McKinley:1948zz,Brown:1970te}, see discussion in Sec.~\ref{subsec:proton_state}.

\paragraph{Unpolarized cross section}

The leading-$Q^2$ behaviour of the inelastic TPE correction to the unpolarized cross section reads \cite{Brown:1970te,Gorchtein:2014hla}
\begin{eqnarray}
\delta_{2\gamma}&=&\frac{Q^2}{2\pi^3}\int_{\bar\nu_\pi}^\infty \frac{\dd \bar \nu}{\bar \nu}\sigma_T (\bar \nu) \ln \left(\frac{4 M^2\,\bar \nu^2}{Q^2\,W^2}\right)\eqlab{delta2galowt} \nonumber\\
&\times&\bigg[\left(1+\frac{\bar \nu^2}{2\epsilon_1^2}\right)\ln \left\vert \frac{\epsilon_1+\bar \nu}{\epsilon_1-\bar \nu} \right\vert\nonumber\\&&+\frac{\bar \nu}{\epsilon_1}\ln \left\vert 1-\frac{\epsilon_1^2}{\bar \nu^2} \right\vert-\frac{\bar \nu}{\epsilon_1}\bigg],
\end{eqnarray}
where $W^2$ is the invariant mass of the hadronic state in the internal Compton scattering process, and $\bar\nu_\pi$ is the pion production threshold. This formula has been first published by Brown \cite{Brown:1970te} in the early 1970's. Later, \Eqref{delta2galowt} has been re-evaluated by Gorchtein \cite{Gorchtein:2014hla} based on the phenomenological fit of the total photoabsorption cross section in Ref.~\cite{Gorchtein:2011xx}. It is important to point out that the coefficient of the order $Q^2 \ln Q^2$ term in $\delta_{2\gamma}$ is a model-independent result.

Tomalak and Vanderhaeghen have extended this approach beyond the leading $Q^2 \ln Q^2$ term and approximate the $\delta_{2\gamma}$ correction for $ep$ scattering in the region up to $Q^2=0.25\,\text{GeV}^2$ in Ref.~\cite{Tomalak:2015aoa} and for 
$\mu p$ scattering in the region up to $Q^2=0.08\,\text{GeV}^2$ in Ref.~\cite{Tomalak:2015hva}, see Fig.~\ref{fig.NearForw}. As can be seen from Eqs.~\ref{eq:T1DR} and \ref{eq:T2DR}, in general, this requires input for the unpolarized proton structure functions, $F_1(x,\bar Q^2)$ and $F_2(x,\bar Q^2)$ (or equivalently, the transverse and longitudinal cross sections, $\sigma_T(\bar \nu,\bar Q^2)$ and $\sigma_L(\bar \nu,\bar Q^2)$), as well as the subtraction function $T_1(0,\bar Q^2)$, where contrary to \Eqref{delta2galowt} the $\bar Q^2$ dependence of the input is kept. Besides the limited applicability range of the near-forward approximation, the $T_1(0,\bar Q^2)$ subtraction function introduces an additional model dependence for the scattering of massive leptons, while its contribution in the elastic $ep$ scattering is suppressed by the electron mass and negligible~\cite{Tomalak:2015aoa}.


\paragraph{Beam normal spin asymmetry}

The leading-$Q^2$ behaviour of the inelastic TPE contribution to the beam normal single spin asymmetry in the diffractive limit (high-energy and forward scattering) reads  \cite{Afanasev:2004pu,Gorchtein:2005za} 
\begin{eqnarray}
\eqlab{BnlowQ}
B_n&=&\frac{m \,Q \,\sigma_T}{8\pi^2}\frac{G_E(Q^2)}{\tau \,G_M^2(Q^2)+\varepsilon \,G_E^2(Q^2)}\nn \\
&\times&\left[2-\ln \frac{Q^2}{m^2}\right],
\end{eqnarray}
where the total photoabsorption cross section is assumed to be independent of the photon virtuality and roughly constant as a function of the invariant mass of the photon-proton system. In the nuclear resonance region, where $\sigma_T$ strongly varies with energy, the cross section enters through an energy-weighted integral. Note that this result is not fully model independent, as it makes use of the Callan-Gross relation \cite{Callan:1969uq} between the longitudinal and the transverse cross section in the DIS region.
The expression \Eqref{BnlowQ} has different dependence on momentum transfer $Q$ and the beam energy $\epsilon_1$ compared to contributions from the proton intermediate state. First, it does not fall off with the beam energy as $m/\epsilon_1$; second, it has linear dependence on $\theta$ at a fixed energy $\epsilon_1$ in the near-forward kinematics, whereas for elastic intermediate states (as well as for a celebrated Mott formula) the asymmetry falls as the third power, $\propto\theta^3$. As a result, the contribution of inelastic intermediate states becomes dominant for energies above GeV \cite{pasquini2004resonance, Afanasev:2004pu, Gorchtein:2005za}.

In Ref.~\cite{Gorchtein:2006mq}, subleading terms in the $Q^2$ expansion of $B_n$ were derived and the energy dependence of the photoabsorption cross section was taken into account. A comparison to \Eqref{BnlowQ} showed that the leading-term in the $Q^2$ expansion is indeed a good approximation. Note that while $B_n$ also has a double-logarithmic enhancements due to hard collinear quasi-real photons \cite{Afanasev:2004pu,Afanasev:2004hp,Gorchtein:2005yz}, it is highly suppressed for small scattering angles as it stems only from the helicity-flip Compton amplitude. However, double-logarithmic enhancement takes place at large scattering angles from a region of quasi-real Compton scattering \cite{pasquini2004resonance}, where both exchanged photons have low virtualities, while the mass of the excited hadronic state is the largest allowed by the incoming beam energy. 

Below the threshold of two-pion production, the absorptive part of proton Compton amplitude that governs the beam asymmetry can be described, as required by unitarity, in terms of quadratic combinations of single-pion production amplitudes. This program was realized in Ref.~\cite{pasquini2004resonance} where, using MAID amplitudes for electroproduction of pions, the authors reached good agreement for the beam asymmetries measured at MAMI for a broad range of scattering angles.



\subsubsection{Medium momentum transfer}\label{MediumMT}


The dispersion relation approach is an appropriate tool to improve our knowledge of inelastic contributions to the most uncertain TPE corrections for arbitrary scattering angles in elastic $ep$ scattering at GeV energies and below. The total TPE correction is given as a sum over all possible intermediate states, proton, pion-nucleon, resonances, etc. As was mentioned in Sec.~\ref{subsec:proton_state}, the elastic (proton) intermediate state was included within unsubtracted dispersion relations in Ref.~\cite{Borisyuk:2006fh} and within subtracted dispersion relations in Refs.~\cite{Tomalak:2014sva,Tomalak:2018jak}. The first inelastic pion-nucleon intermediate state was evaluated from data-driven MAID parameterization for the pion electroproduction amplitudes~\cite{Drechsel:1998hk,Drechsel:2007if} at low momentum transfers ($Q^2 \lesssim 0.064~\mathrm{GeV}^2$) in Ref.~\cite{Tomalak:2016vbf}. This straightforward, but numerically involved calculation, consistently accounts for the contribution of all resonances and non-resonant background. Going to higher momentum transfers requires analytical continuation of $ep$ scattering amplitudes to the unphysical region of kinematics. Such a novel method of analytical continuation for multiparticle intermediate states was developed and tested in Ref.~\cite{Tomalak:2017shs}. It allowed authors to evaluate pion-nucleon contributions to TPE corrections in elastic $ep$ scattering at more than an order of magnitude larger momentum transfers ($Q^2 \lesssim 1~\mathrm{GeV}^2$) compared to their previous work~\cite{Tomalak:2016vbf}.

For higher $Q^2$, the dispersive approach for resonant intermediate states was developed in Ref.~\cite{Blunden:2017nby}, and implemented for spin-1/2 and spin-3/2 resonances with mass below 1.8~GeV in Ref.~\cite{Ahmed:2020uso}. They included a Breit-Wigner shape with a finite width for each resonance, and for the resonance electrocouplings at the hadronic vertices they used helicity amplitudes extracted from the analysis of CLAS electroproduction data~\cite{HillerBlin:2019jgp}.

\begin{figure}[ht]
\includegraphics[width=0.45\textwidth]{TPE/fig/dSig.pdf}
\caption{The TPE correction $\delta_{2\gamma}$ (in \%) versus $\varepsilon$ for nucleon plus all
spin-parity $1/2^\pm$ and $3/2^\pm$ states \cite{Ahmed:2020uso} at
    $Q^2 = 0.2$~GeV$^2$ (green dashed line),
    0.5~GeV$^2$ (dark red long-dashed),
    1~GeV$^2$ (red solid),
    3~GeV$^2$ (orange dot-dashed),
    4~GeV$^2$ (black dotted), and
    5~GeV$^2$ (blue dashed). The shaded bands correspond to
    the uncertainty propagated from the input electrocouplings.}
\label{fig.Dsig}
\end{figure}

The combined effect on $\delta_{2\gamma}$ from the nucleon plus all the spin-parity $1/2^\pm$ and $3/2^\pm$ resonances is illustrated in Fig.~\ref{fig.Dsig} as a function of $\varepsilon$ for a range of fixed $Q^2$ values between 0.2 and 5~GeV$^2$.
At low $Q^2\lesssim 3$~GeV$^2$ the net excited state resonance contributions are small, and the total correction is dominated by the nucleon elastic intermediate state.
The net effect of the higher mass resonances is to increase the magnitude of the TPE correction at $Q^2 \gtrsim 3$~GeV$^2$, due primarily to the growth of the (negative) odd-parity $N(1520)~\!3/2^-$ and $N(1535)~\!1/2^-$ resonances, which overcompensates the (positive) contributions from the $\Delta(1232)~\!3/2^+$. 
At the highest $Q^2=5$~GeV$^2$ value, the total TPE correction $\delta_{\rm tot}$ reaches $\approx 6\% - 7\% $ at low $\varepsilon$.

An estimate of the theoretical uncertainties on the TPE contributions can be made by propagating the uncertainties on the fitted values of the transition electrocouplings~\cite{HillerBlin:2019jgp}, which are dominated by the $\Delta(1232)~3/2^+$ and $N(1520)~3/2^-$ intermediate states.
At low~$Q^2$, $Q^2 \lesssim 0.5$~GeV$^2$, the uncertainties are insignificant, but become more visible at higher $Q^2$ values, as illustrated by the shaded bands in Fig.~\ref{fig.Dsig}, for $Q^2 = 1-5~ \rm{GeV}^2$.



\subsubsection{Large momentum transfer}

For TPE calculations at large momentum transfers, quark degrees of freedom have to be included. In an approach based on generalized parton distributions and a handbag mechanism (with two photons coupling to the same quark) \cite{Chen:2004tw,Afanasev:2005prd}, TPE effects for the unpolarized cross section and for polarization asymmetries were calculated. Perturbative QCD (pQCD) approaches were considered in Refs.~\cite{Borisyuk:2008db, Guttmann:2010au}. In particular, in Ref.~\cite{Guttmann:2010au}, special consideration was given to the coupling of two photons to different quarks, noticing that a large momentum may be transferred with one less gluon compared to the case of pQCD description of elastic proton FFs.


\subsection{Effective field theory calculations}\seclab{EFTCalc}

For elastic $\ell p$ scattering at lepton energies well below the proton mass, we can use low-energy effective field theories (EFTs). At such energies, we do not need the full knowledge of the non-perturbative properties of the proton. For OPE, we only need a few effective couplings, e.g., the proton charge, magnetic moment, and charge radius. Similar simplifications apply to TPE. In the following, we discuss two such effective theories: heavy baryon chiral perturbation theory (HBChPT) and QED-NRQED, and review the current status of their predictions for TPE.
Before a detailed discussion, it is worth mentioning that radiative corrections in the elastic electron-proton scattering when $Q^2 >> m^2_e$ formulated recently in the Soft-Collinear Effective Theory (SCET) approach within QED~\cite{Hill:2016gdf}, which allows us also for a systematic resummation of large logarithms. Moreover, the two-photon exchange corrections in the hard momentum region were also evaluated in the SCET framework within QCD in Ref.~\cite{Kivel:2012vs}.

\subsubsection{Heavy baryon chiral perturbation theory}

Within this topical collection, an  exact analytical evaluation of the TPE contribution to elastic $\ell p$ scattering in low-energy EFT, namely, HBChPT  at NLO, is published \cite{choudhary2023analytical}.
The LO HBChPT
Lagrangian contains one derivative only. The NLO Lagrangian  will have two derivatives and terms of order $1/M$, where $M$ is the proton mass. 
In this framework, the proton being the heavy degree of freedom is treated non-relativistically. The leptons, the electron and muon, along with the photons constituting the light degrees of freedom, are treated in standard covariant QED. Throughout, the masses of the relativistic leptons are kept non-zero. 

Previously, the NLO HBChPT prediction had been calculated using the soft-photon approximation (SPA) \cite{Talukdar:2019dko,Talukdar:2020aui}. As before, Ref.~\cite{choudhary2023analytical} assumes 
that at low energies the dominant photon loop contributions arise only from the elastic proton intermediate state.  The inelastic proton intermediate states are beyond the intended accuracy of the evaluation. Notably, since the proton FFs, which parametrize the hadronic structure effects, enter in HBChPT at NNLO, the proton is point-like at NLO. All hard- and soft-photon exchanges are now included in the two-photon loop corrections which are evaluated exactly. The methodology relies on the successive use of partial fractions and integration by parts facilitating the decomposition of the rather intricate TPE four-point loop functions into a system of two- and three-point scalar master integrals, which are straightforward to evaluate analytically. 
 
To LO accuracy, our low-energy TPE result approximates the well-known McKinley-Feshbach two-photon correction in potential scattering theory, cf.\ \Eqref{Feshbach}. At NLO the finite analytical real parts of the TPE loop contributions to the elastic differential cross section are presented. The exact analytical evaluation of the NLO bremsstrahlung process is not included. This means that the systematic cancellations of the IR-singular terms that arise from the TPE loops are not discussed. While the details of the TPE evaluation are going to be presented in Ref.~\cite{choudhary2023analytical}, 
a rigorous NLO evaluation of the bremsstrahlung subtracted IR-free TPE contribution to the radiative corrections shall be a subject matter of a later publication.

\subsubsection{QED-NRQED}

An EFT related to HBChPT, is QED-NRQED, suggested in \cite{Caswell:1985ui,Pineda:1997bj,Pineda:1998kj,Hill:2012rh} and further developed in \cite{Dye:2016uep,Dye:2018rgg,Peset:2014jxa}. In this EFT the proton is treated non-relativistically, using NRQED, and the leptons and photons are treated using QED. Unlike HBChPT, the EFT does not include other hadronic degrees of freedom, such as pions. Such effects are encoded in the QED-NRQED coupling constants. Thus QED-NRQED is most useful for elastic $\ell p$ scattering.  

Up to dimension six, the independent NRQED Wilson coefficients are $c_F$ and $c_D$ corresponding to the proton magnetic moment and charge radius, respectively. There are also two dimension-six contact interaction terms: a spin-independent term and a spin-dependent term \cite{Hill:2012rh}, 
\begin{equation}
{\cal L}_{\ell\psi}^{d=6}=\dfrac{b_1}{M^2}\psi^\dagger\psi\,\bar \ell\gamma^0\ell+\dfrac{b_2}{M^2}\psi^\dagger\sigma^i\psi\,\bar \ell\gamma^i\gamma^5\ell\,,
%{\cal O}\left(1/M^3\right)
\end{equation}
where $\psi$ ($\ell$) is the proton (lepton) field and $M$ is the proton mass, the cutoff scale of the EFT. 

In Ref.~\cite{Dye:2016uep}, it was shown that QED-NRQED $\ell p$ scattering at ${\cal O}(Z\alpha)$ and power $1/M^2$  reproduces the known Rosenbluth scattering formula expanded to $1/M^2$. It requires just the Dirac Lagrangian and the NRQED Lagrangian up to $1/M^2$, implying that $b_1$ and $b_2$ are zero at ${\cal O}(Z\alpha)$. QED-NRQED $\ell p$ scattering at ${\cal O}(Z^2\alpha^2)$ and leading power in $1/M$ reproduces the ${\cal O}(Z^2\alpha^2)$ terms in the scattering of a lepton off a static $1/r$ potential \cite{Dalitz:1951ah}. Here $Z$ denotes the proton charge in units of $|e|$ and is used to denote the appropriate class of radiative corrections.  



Since the Wilson coefficients $b_1$ and $b_2$ are zero at ${\cal O}(Z\alpha)$, they are sensitive to TPE at scales above the proton mass. In Ref.~\cite{Dye:2018rgg}, $b_1$ and $b_2$ were calculated at ${\cal O}(Z^2\alpha^2)$. They were extracted by calculating the $\ell p\to\ell p$ off-shell  forward scattering amplitude at ${\cal O}(Z^2\alpha^2)$ and power $1/M^2$ in the effective and full theory  in both Feynman and Coulomb gauges. Two cases were considered:  a toy example of a non-relativistic point particle (p.p.), and the real proton described by a hadronic tensor, namely, the forward VVCS off the proton discussed in \secref{VVCS}.


For the toy example, Ref.~\cite{Dye:2018rgg} found $b_1^{\text{ p.p.}}=0$ and $b_2^\text{ p.p.}=Q_\ell^2Z^2\alpha^2\left[{16}/{3}+\ln\left({M}/{2\Lambda}\right)\right]$, where $\Lambda$ is the UV cutoff of QED-NRQED, and $Q_\ell$ is the lepton charge in units of $|e|$. Surprisingly, $b_1^{\text{ p.p.}}=0$ at ${\cal O}(Z^2\alpha^2)$. For the case of the real proton, Ref.~\cite{Dye:2018rgg} derived implicit expressions for $b_1$ and $b_2$ in terms of the components of the hadronic tensor. Considering only the contribution to the Wilson coefficients of $F_1(0)$, $F_2(0)$ and $M^2 F_1^\prime(0)$, related to the proton charge, magnetic moment and charge radius respectively, Ref.~\cite{Dye:2018rgg} finds: 
\begin{align}
b_1/(\alpha^2Q_\ell^2)&=0+\cdots\,,\\
\nn\\
b_2/(\alpha^2Q_\ell^2)&=F_1(0)^2\left[\frac{16}{3}+\ln\left(\frac{M}{2\Lambda}\right)\right] \nn\\&+\frac{16}{3}F_1(0)F_2(0)\nn\\&+\frac{F_2(0)^2}{2}\left[\frac{17}{12}-\ln\left(\frac{M}{2\Lambda}\right)+3\ln\left(\frac{Q}{M}\right)\right]\nn\\&+\cdots\,. 
\end{align}
The ellipsis denotes  $\mbox{non } F_1(0), F_2(0), M^2 F_1^\prime(0)$ terms.  Surprisingly, again there is no contribution to $b_1$. It does not follow from a symmetry of the EFT and it might be a one-loop ``accident". %See the talk by Gil Paz in these proceedings for  details. 


This implies that low-energy elastic $\ell p$ scattering is much less sensitive to spin-independent TPE effects above the proton mass scale compared to spin-dependent ones. On the other hand, the proton charge radius extraction will be more robust. 

With the calculation of the QED-NRQED dimension six couplings: $c_F,c_D, b_1, b_2$, at hand, the next logical step is the calculation of the differential cross section for elastic $\ell p$ scattering. Such a calculation was not performed in the literature yet. Once performed, it would be interesting to compare its result to HBChPT. 



\subsection{Empirical extractions of TPE}\seclab{TPEempirical}

Besides theoretical predictions, it is possible to extract TPE amplitudes, as well as the TPE contributions to the unpolarized cross section $\delta_{2\gamma}$,  from empirical cross section and polarization transfer measurements.

\subsubsection{Phenomenological fits of \texorpdfstring{$\delta_{2\gamma}$}{TEXT}}

The size of the TPE contribution becomes more significant at larger $Q^2$. A phenomenological fit of the $Q^2$-dependent part, as an addition to the Feshbach correction, \Eqref{Feshbach}, is \cite{A1:2013fsc} 
\begin{equation}
    \delta_{2 \gamma} =-(1-\varepsilon) \, a \ln (b\, Q^2+1),
\end{equation}
where $a=0.069, \,b=0.394\, \text{GeV}^{-2}$. This fit includes the world data set on unpolarized cross section measurements (updated to a common standard on radiative corrections) and measurements of the FF ratio using polarization and assumes that the whole discrepancy is driven by hard TPE.

Similarly, Schmidt in Ref.~\cite{Schmidt:2019vpr} uses several existing parametrizations for FFs from Rosenbluth-type experiments and polarization measurements of the ratio to determine the hard TPE contribution, under the assumption that TPE preserves the linearity of the reduced cross section in $\varepsilon$.

\subsubsection{Updated extraction of TPE amplitudes}

A measurement of the three observables for elastic $ep$ scattering, $\sigma_R, P_l,$ and $P_t$,  as a function of $\varepsilon$ for a fixed value of $Q^2$ allows one to extract in a model-independent fashion the three TPE amplitudes, which conserve the lepton helicity, as can be seen from Eqs.~\ref{crossen}, \ref{PtPlradcorr} and \ref{pl_2gamma}. Such an analysis has been performed in \cite{Guttmann:2010au} at $Q^2 = 2.5$~GeV$^2$, based on 
available cross section data as well as the $\varepsilon$-dependence of  
  $P_t/P_l$ and $P_l/P_l^{Born}$ 
  as measured by the JLab/Hall C experiment~\cite{GEp2gamma:2010gvp}. 
    The combination of both unpolarized and polarization experiments at a same $Q^2$ value 
provided the necessary three observables to extract the 
  $\varepsilon$-dependence of the three TPE amplitudes $Y_M$, $Y_E$, and $Y_3$, introduced in \Eqref{2gaampl}. A concise update of this analysis in view of more recent data is given here. 

The data for $P_t / P_l$ of the dedicated JLab/Hall C GEp-$2 \gamma$ 
experiment~\cite{GEp2gamma:2010gvp,Puckett:2017flj}, as shown in the upper panel of Fig.~\ref{fig:PtPl},  
do not see any systematic TPE effect within their error bars of order 1\%.  
One can therefore effectively fit the observable on the lhs of \Eqref{Rratios}   
assuming an $\varepsilon$-independent part, which equals its OPE limit,
\begin{eqnarray}
- \sqrt{ \frac{ \tau (1 +\varepsilon)}{2\varepsilon}}  \frac{P_t}{P_l} \simeq 						
	R_{EM}(Q^2),
\end{eqnarray}
and extract the $Q^2$-dependence of $R_{EM}$, cf.\ \Eqref{REM}, from this observable, using the parameterization given in Ref.~\cite{Puckett:2017flj} (Global fit II).


\begin{figure}[ht]
\includegraphics[width=0.45\textwidth]{TPE/fig/R.pdf}
\includegraphics[width=0.45
\textwidth]{TPE/fig/Pl2Plborn.pdf}
\caption[]{Polarization transfer observables for elastic $ep$ scattering at $Q^2 = 2.5$~GeV$^2$. Upper panel: $- \mu_p \sqrt{\tau (1 + \varepsilon)/(2 \varepsilon)} P_t / P_l$, lower panel: $P_l / P_l^\text{Born}$. The JLab/Hall C data are from the GEp-$2 \gamma$ experiment~\cite{GEp2gamma:2010gvp}, 
using the updated analysis of \cite{Puckett:2017flj}. The red curves are the fit to the data. Left: $\varepsilon$-independent fit according to 
Ref.~\cite{Puckett:2017flj}; right: updated fit from Ref.~\cite{Guttmann:2010au} according to Eq.~(\ref{eq:plfit}). 
Note that the values for the polarization observables are shown without radiative corrections. 
The bulk of these corrections (virtual corrections on the lepton side and soft-photon emission  corrections), which factorize in terms of the Born cross section, drop out of the asymmetries, and the hard TPE is therefore expected to be the leading correction. It has been checked in Refs.~\cite{GEp2gamma:2010gvp,Puckett:2017flj} using the MASCARAD program~\cite{Afanasev:2001jd} that the standard radiative corrections yield a multiplicative correction around or less than 0.1\% on the asymmetries.}
\label{fig:PtPl}
\end{figure}

For the longitudinal polarization observable $P_l/ P_l^\text{Born}$, shown in the lower panel of Fig.~\ref{fig:PtPl}, the deviation from unity is 6.2 times the statistical uncertainty and 2.2 times the total uncertainty~\cite{Puckett:2017flj}. The TPE correction effect has been parameterized as~\cite{Guttmann:2010au}
\begin{eqnarray}
P_l / P_l^\text{Born} &=& 1 + a_l \, \varepsilon^4 (1 - \varepsilon)^{1/2}.
\label{eq:plfit}
\end{eqnarray}
The comparison of this fit function to the data is shown in the lower panel of Fig.~\ref{fig:PtPl} for  
$a_l = 0.09$.  

Recently, an updated global analysis for $\sigma_R$, including new high-$Q^2$ data from JLab, has been given in Ref.~\cite{Christy:2021snt}.
In Ref.~\cite{Christy:2021snt}, the data for the reduced cross section have been fitted by a linear in $\varepsilon$ behavior
\begin{eqnarray}
\frac{\sigma_R(\varepsilon, Q^2)}{(\mu_p G_D)^2} = a + \varepsilon b,   
\label{eq:sigmaRfit}
 \end{eqnarray}
 with $\mu_p$ the proton magnetic moment, and $G_D(Q^2) \equiv 1/(1 + Q^2/0.71)^2$ the standard dipole. 
 As the TPE amplitudes vanish for $\varepsilon \to 1$, the proton FF $G_M$ can be extracted from Eqs.~\ref{eq:sigmaR} and \ref{eq:sigmaRfit} as
\begin{eqnarray}
 \left( \frac{G_{Mp} }{ \mu_p G_D} \right)^2 = \frac{a + b }{1 + \frac{1}{\tau} R_{EM}^2}.
 \label{eq:GMfit}
 \end{eqnarray}
The parameterization for $\sigma_R$ chosen in~\cite{Christy:2021snt} does however not have the correct $\varepsilon \to 1$ limit, in which  
the TPE amplitudes vanish.  
At a fixed value of $Q^2$, the limit $\varepsilon \to 1$ corresponds with the Regge (high-energy) limit, as 
$Q/(2\eps_1) \simeq (1 - \varepsilon)^{1/2} / \sqrt{2}$ for $\varepsilon \to 1$, with $\eps_1$ the electron lab energy.
In an expansion around $\varepsilon \to 1$, the leading TPE correction can be model independently expressed as~\cite{Tomalak:2015aoa} 
\begin{align}
& \delta_{2 \gamma}(\varepsilon \to 1, Q^2) = \alpha \pi \left(\frac{Q}{2 \eps_1}\right) + \frac{\alpha }{\pi} \left(\frac{Q^2}{M\eps_1} \right)  \ln^2 \left(\frac{Q}{2\eps_1} \right) \nn \\ &
+ c \left(\frac{Q^2}{M\eps_1} \right)  \ln \left(\frac{Q}{2\eps_1} \right) + 
 {\mathcal O}\left( \frac{Q^2}{M^2}, \frac{Q^2}{M \eps_1}, \frac{Q^2}{\eps_1^2}\right).
 \label{eq:expansion}
 \end{align}
In Eq.~\ref{eq:expansion}, the first term, which dominates for $\varepsilon \to 1$ corresponds with the McKinley-Feshbach TPE correction due to the scattering of relativistic electrons in the Coulomb field of the proton~\cite{McKinley:1948zz}, the $Q^2 \ln^2 Q$ term is due to the elastic intermediate state only, and the $Q^2 \ln Q$ term contains the effect due to inelastic intermediate states~\cite{Brown:1970te,Tomalak:2015aoa}, cf.\ discussion in Sec.~\ref{SmallMT}.



In order to combine the empirical observation of a linearity of $\sigma_R$ over most of the $\varepsilon$ region with the correct $\varepsilon \to 1$ limit, the fit function of Eq.~\ref{eq:sigmaRfit} for $\sigma_R$ can be improved as
\begin{align}
&\frac{\sigma_R(\varepsilon, Q^2)}{(\mu_p G_D)^2} = a + \varepsilon b \nn \\
&+ (a + b) \frac{ \left(1 + \frac{\varepsilon}{\tau}  R_{EM}^2 \right) }{ \left(1 + \frac{1}{\tau}  R_{EM}^2 \right) } \frac{\alpha \pi}{\sqrt{2}} \left[(1 - \varepsilon)^{1/2} - (1 - \varepsilon)\right],   
\label{eq:sigmaRfit2}
 \end{align}
 which does not change the fit value for small $\varepsilon$
 \begin{eqnarray}
\frac{\sigma_R(\varepsilon = 0, Q^2)}{(\mu_p G_D)^2}  = a, 
\end{eqnarray}
and has the correct limit for $\varepsilon \to 1$:
\begin{eqnarray}
\frac{\sigma_R(\varepsilon \to 1, Q^2)}{(\mu_p G_D)^2} &=& (a + b) \left[ 1 + \frac{\alpha \pi}{\sqrt{2}} (1 - \varepsilon)^{1/2}  \right]. \qquad
 \end{eqnarray}
Using the fit function of Eq.~\ref{eq:sigmaRfit2}, the TPE correction $\delta_{2 \gamma}$ follows as
\begin{align}
\delta_{2 \gamma}(\varepsilon, &Q^2) = \frac{\alpha \pi}{\sqrt{2}} (1 - \varepsilon)^{1/2} \nn \\
&+ (1 - \varepsilon) \left[ 
\frac{\left( - b + \frac{a}{\tau} R_{EM}^2 \right) }{(a + b) \left( 1 + \frac{\varepsilon}{\tau} R_{EM}^2 \right)} - \frac{\alpha \pi}{\sqrt{2}}\right] .
\label{eq:delta2gafin}
\end{align}
In Fig.~\ref{fig:sigmaR} (upper panel), we compare, for $Q^2$ = 2.5 GeV$^2$, the fit functions for $\sigma_R$ of Ref.~\cite{Christy:2021snt} (red dotted curve) with the modified fit function of Eq.~\ref{eq:sigmaRfit2} (green solid curve), as well as $\sigma_R$ in the OPE approximation, setting $\delta_{2 \gamma} = 0$, (blue dashed curve).
The resulting $\varepsilon$-dependence of the TPE correction $\delta_{2 \gamma}$ for $Q^2$ = 2.5 GeV$^2$, based on the linear fit of Ref.~\cite{Christy:2021snt}, using the improved fit function of Eq.~\ref{eq:sigmaRfit2}, using the spline and Pad\'e fits of Ref.~\cite{A1:2013fsc}, are compared in Fig.~\ref{fig:sigmaR} (lower panel). For comparison, the McKinley-Feshbach correction, which gives the leading behavior for $\varepsilon \to 1$, is also shown (black dot-dashed curve).  

\begin{figure}[ht]
\includegraphics[width=0.45\textwidth]{TPE/fig/sigmaR.pdf}
\includegraphics[width=0.45\textwidth]{TPE/fig/Delta2Gamma.pdf}
\caption[]{Upper panel: $\varepsilon$-dependence of the reduced cross section $\sigma_R$ for $Q^2$ = 2.5 GeV$^2$, based on the linear fit of Ref.~\cite{Christy:2021snt} (red dotted curve), using the modified fit function of Eq.~\ref{eq:sigmaRfit2} (green solid curve), as well as in the OPE approximation, setting $\delta_{2 \gamma} = 0$, (blue dashed curve). Lower panel: $\varepsilon$-dependence of the TPE  correction $\delta_{2 \gamma}$ for $Q^2$ = 2.5 GeV$^2$, based on the linear fit of Ref.~\cite{Christy:2021snt} (red dotted curve), using the modified fit function of Eq.~\ref{eq:sigmaRfit2} (green solid curve), using the spline fit (purple long-dashed curve) and Pad\'e fit (yellow short-dashed curve) of Ref.~\cite{A1:2013fsc}. The black dot-dashed curve shows the McKinley-Feshbach correction, which gives the leading behavior for $\varepsilon \to 1$. }
\label{fig:sigmaR}
\end{figure}

\begin{figure}[ht]
\includegraphics[width=0.45\textwidth]{TPE/fig/YME3.pdf}
\caption[]{$\varepsilon$ dependence for the TPE amplitudes for $Q^2 = 2.5$~GeV$^2$, based on the polarization data of 
Ref.~\cite{Puckett:2017flj}, shown in Fig.~\ref{fig:PtPl}, and the cross section fit of Ref.~\cite{Christy:2021snt}, improved using the fit function of Eq.~\ref{eq:sigmaRfit2}, given by the red curves in Fig.~\ref{fig:sigmaR}.}
\label{fig:YME3}
\end{figure}

Based on the measured $\varepsilon$-dependence 
of the three elastic $e^-p$ observables ($\sigma_R,\, P_l$, and $P_t$ at $Q^2= 2.5$~GeV$^2$ as shown in Figs.~\ref{fig:PtPl} and \ref{fig:sigmaR}), an updated empirical extraction 
of the three TPE amplitudes, 
$Y_M, Y_E$, and $Y_3$, at this $Q^2$ value was obtained, see Fig.~\ref{fig:YME3}. 

The updated empirical extraction of the TPE correction $\delta_{2\gamma}$ for general values of $Q^2$ and $\epsilon$, is compared to data from the OLYMPUS, VEPP-3 and CLAS experiments in Figs.~\ref{fig.OLYMPUS}, \ref{fig.VEPP} and \ref{fig.CLAS} (dot-dashed pink curves), respectively. Note that the empirical extraction displays a small oscillation in the low-$Q$ region, cf.\ $\varepsilon=0.45, 0.88$ panels in Fig.~\ref{fig.CLAS}. This, however, is not a true feature but a remnant that propagated from the underlying fit of $R_{EM}(Q^2)$ that was used~\cite{Puckett:2017flj}, and is enhanced due to two opposite sign contributions in the $(1 - \varepsilon)$ term in Eq.~\ref{eq:delta2gafin}.




\subsection{Dedicated experiments and comparison to theory}\label{Sec:Experiments}

\subsubsection{Single-spin asymmetries}\seclab{BeamAsymmetry}

As stated above in Sec.~\ref{SSAintro}, parity-conserving SSA for elastic $ep$ scattering are zero in the first Born approximation. Therefore measurements of either target or beam SSA may provide useful information on TPE amplitudes. While beam asymmetries arising from TPE are of the order $10^{-5}$ for electrons \cite{Afanasev:2004pu,Gorchtein:2004ac, pasquini2004resonance} and $10^{-3}$ for muons \cite{Koshchii:2019b}, the target SSA are expected at percent level \cite{Afanasev:2005prd}. The beam asymmetries were measured with high accuracy at MIT \cite{Wells:2001bn}, MAMI \cite{maas2005measurement} and Jefferson Lab \cite{armstrong2007transverse,Androic:2012bn,Abrahamyan:2012bn,PREX:2021uwt} using experimental setups designed to study parity-violating electron-helicity asymmetries. Measurements of target SSA were performed at Jefferson Lab on a polarized $^3$He target \cite{Zhang:2015kna}; results for $Q^2>$ 1 GeV$^2$ were in agreement with GPD-based calculations \cite{Afanasev:2005prd} for a neutron. 

The measurements of SSA in elastic electron scattering provided unambiguous evidence of TPE effects, while probing contributions arising from an absorptive part of the hadronic Compton amplitude entering the TPE mechanism. 



\subsubsection{Cross sections and charge asymmetries}

The TPE correction, defined in Eq.~\ref{delta2gTPE}, can be accessed experimentally by measuring the ratio of cross sections from elastic $\ell ^+p$ vs. $\ell^- p$ scattering. Corrections that depend on an odd power of the lepton charge do not cancel in this ratio. This includes the leading TPE correction, $\delta_{2\gamma}$. It also includes, however, the interference of lepton and proton bremsstrahlung radiation,  $\delta_{b}$, which is comparable to $\delta_{2\gamma}$. Their sum gives us the ``odd'' part of the radiative corrections, $\delta_{\text{odd}}=\delta_{b}+\delta_{2\gamma}$.
By contrast, the ``even'' part of the corrections, $\delta_{\text{even}}$, include logarithmically enchanced terms, $\log (Q^2/m_{\ell})$, coming from lepton bremsstrahlung, vacuum polarization and vertex corrections. As a result, the even part is relatively large, implying $|\delta_{\text{odd}}| \ll |\delta_{\text{even}}|$.

The above mentioned cross section ratio can be written in the following way \cite{Afanasev:2017gsk}:
\begin{align}
    R_{\pm}^\text{exp.}&=\frac{\sigma(\ell ^+p)}{\sigma(\ell^-p)} = \frac{1+\delta_{\text{even}}-\delta_{\text{odd}}}{1+\delta_{\text{even}}+\delta_{\text{odd}}}\nn\\& \approx 1-\frac{2 \,\delta_{\text{odd}}}{1+\delta_{\text{even}}}.
\end{align}
The TPE correction can be extracted from 
\begin{equation}
\eqlab{R2gammadef}
       R_{2\gamma} \approx 1-2 \,\delta_{2\gamma},
\end{equation}
which is the 
measured ratio $R_{\pm}^\text{exp.}$ corrected for  $\delta_{b}$ and $\delta_{\text{even}}$.


Experimental possibilities are limited, as such a measurements requires both electron and positron beams of high quality and at relevant beam energies of a few GeV. In recent years, three experiments have published results, cf.\ Figs.~\ref{fig.OLYMPUS}-\ref{fig.CLAS}: an experiment \cite{Rachek:2014fam} at the VEPP-3 storage ring measured the cross section ratio with beam energies of 1.0 and 1.6 GeV with a non-magnetic spectrometer covering angles between 15$^\circ$ and 105$^\circ$. Since the experiment lacks a relative luminosity determination of sufficient precision, the data is published relative to the most forward point set to a ratio of one. However, this point is sufficiently far away from $\ep=1$ that a proper comparison with curves requires a shift of the data to match the predictions at this normalization point. CLAS at Jefferson Lab \cite{CLAS:2013mza,CLAS:2014xso} generated a wide-energy-spread electron and positron beam by converting the CEBAF electron beam to a photon beam via a converter target, subsequent pair-production from that photon beam, and by cleaning and recombining the electron/positron beam via a chicane magnet system. Using the detector setup of the CLAS collaboration in Hall B, the so produced wide-energy beam-proton scattering events are detected. By the curvature of the scattered lepton in the magnetic field, the charge of the particle can be determined.  Because of the wide-energy spread of the beam, the resulting data is sorted in large bins in $\ep$ and $Q^2$. The collaboration published data with several (non-independent) bin selections. The OLYMPUS experiment \cite{Henderson:2016dea} at the DORIS ring at DESY, Hamburg, measured using a monoenergetic beam with an energy of around 2 GeV using an open-cell internal hydrogen target. Scattered leptons and protons were detected with a spectrometer upgraded from the former BLAST detector. OLYMPUS reached the highest $Q^2$ of the three experiments of slightly more than 2~$\GeV^2$.

\subsubsection{Comparison theory vs.\ experiment} 

 The OLYMPUS data (Fig.~\ref{fig.OLYMPUS}) at smaller $Q^2$ / larger $\ep$ are below unity. This is only replicated by the (mostly) phenomenological models, while theoretical predictions produce a ratio below 1 only much closer to $\ep=1$. Overall, the phenomenological predictions compare well to the data, but both predict a somewhat larger effect at larger $Q^2$\,/\,smaller $\ep$, albeit within the uncertainty of the data. Note that, calculations prior to and motivating the experimental proposal  \cite{Guttmann:2010au} predicted a large and visible deviation of the ratio from unity already around $Q^2\simeq$ 2 GeV$^2$.

\begin{figure}[ht] 
\includegraphics[width=\columnwidth]{TPE/fig/OLYMPUS.pdf} 
\caption{Comparison of predictions for $R_{2\gamma}$, \Eqref{R2gammadef}, from $ep$ scattering with beam energy $E_\text{beam}=2.01$ GeV to results from the OLYMPUS experiment (black points; statistical and uncorrelated systematic uncertainties added in quadrature; gray errorband for correlated systematic uncertainties) \cite{Henderson:2016dea}. Theoretical predictions: sum of nucleon plus spin-parity $1/2^\pm$ and $3/2^\pm$ resonances by Ahmed et al.~(dashed orange) \cite{Ahmed:2020uso}, nucleon plus $\pi N$ intermediate state by Tomalak et al.~(dotted blue) \cite{Tomalak:2017shs}. Phenomenological predictions: Bernauer et al.~(solid violet line with error band) \cite{A1:2013fsc}, Vanderhaeghen et al.~(dot-dashed pink) \cite{Guttmann:2010au}.}
\label{fig.OLYMPUS}
\end{figure} 

The VEPP-3 (Fig.~\ref{fig.VEPP}) and CLAS data (Fig.~\ref{fig.CLAS}) also prefer the phenomenological extractions; however, the CLAS data are mostly compatible with Refs.~\cite{Ahmed:2020uso} and \cite{Tomalak:2017shs} as well. The CLAS $\ep=0.88$ data set and the VEPP-3 $E_\text{beam}=1.594$ GeV data set show a ratio below 1 for small $Q^2$ / large $\varepsilon$, similar to the behavior seen in the OLYMPUS result.




\begin{figure}[ht]
\includegraphics[width=\columnwidth]{TPE/fig/VEPPIII.pdf}
\caption{Comparison of predictions for $R_{2\gamma}$, \Eqref{R2gammadef}, from $ep$ scattering with beam energies $E_\text{beam}=0.998$ GeV and $1.594$ GeV to results from the VEPP-3 experiment (black triangles) \cite{Rachek:2014fam} Note that the points are normalized according to the fitting model by Bernauer et al. \cite{A1:2013fsc} at the points $\varepsilon =0.931, \,\, Q^2 =0.128$ $\text{GeV}^2$ (top panel) and $\varepsilon =0.98, \,\, Q^2 =0.097$  $\text{GeV}^2$ (bottom panel). For further details see legend in Fig.~\ref{fig.OLYMPUS}.}
\label{fig.VEPP}
\end{figure}

\begin{figure}[t!]
\includegraphics[width=\columnwidth]{TPE/fig/CLAS.pdf}
\caption{Comparison of predictions for $R_{2\gamma}$, \Eqref{R2gammadef}, from $ep$ scattering with momentum transfer $Q^2=0.85$ GeV,  $1.45$ GeV, and photon polarization parameter $\varepsilon=0.45$ and  $0.88$ to results from the CLAS experiment (black squares) \cite{Rimal:2016toz}. For further details see legend in Fig.~\ref{fig.OLYMPUS}.} 
\label{fig.CLAS}
\end{figure} 






A global comparison of the data from the three experiments and the calculation from  \cite{Kuraev:2007dn} was done in Ref. \cite{Bytev:2019rdc}. 
The largest contribution to the charge asymmetry is due to the soft photon emission. A specific procedure was suggested to limit the effect of the different corrections applied to the data. The results do not favor an enhanced contribution of TPE in the considered kinematical range.



\subsection{Outlook} \seclab{TPEOutlook}

The OLYMPUS data gives a hint that at higher $Q^2$, TPE might not explain all of the discrepancy. Theoretical descriptions, which do not show a particular good agreement in the measured $Q^2$ range, are not at all tested at these larger $Q^2$. Precision measurements of the FFs at these large $Q^2$ however depend on an accurate description of TPE. This makes experimental determinations of TPE at these kinematics highly desirable. Unfortunately, experimental activities are limited by the availability of suitable beams.

The TPEX experiment \cite{Alarcon:2023kuv} seeks to measure TPE at DESY using beams of 2 and 3 \GeV~and up to a $Q^2$ of 4.6 $(\GeV/c)^2$. This measurement would require a new extracted-beam line at the DESY test beam facility. Current plans for future construction at DESY make it unlikely that this experiment can be realized.

Jefferson Lab is planning to construct a positron source for CEBAF. Current timelines put the availability of beam beyond 2030, if funding can be secured. Such a capability would allow for a comprehensive TPE program, with possible measurements in Halls A,~B and C \cite{Arrington:2021kdp,Bernauer:2021vbn,Cline:2021hkr,Grauvogel:2021btg}.

At smaller $Q^2$, MUSE \cite{MUSE:2017dod,Cline:2021ehf} and AMBER \cite{Adams:2018pwt} will take data for $\mu p$ scattering with both beam charges. From this data, TPE can be determined at small $Q^2$ close to the validity of the Feshbach limit, and in the case of AMBER, at $\ep$ extremely close to 1. Theoretical predictions for the inelastic and total TPE in the kinematic region envisaged by the MUSE experiment are shown in Fig.~\ref{fig.NearForw}.

\begin{figure}[ht]
\includegraphics[width=0.45\textwidth]{TPE/fig/NearForwBoth.pdf}
\caption{Predictions for the inelastic (top panel) and total (bottom panel) TPE correction to $\mu^- p$ elastic scattering \cite{Tomalak:2015hva} for the three different beam momenta envisaged by the MUSE experiment \cite{MUSE:2013uhu}. While the elastic TPE dominates, the main uncertainty stems from the inelastic part. The relative uncertainty for the total TPE is thus small and not displayed here.} 
\label{fig.NearForw}
\end{figure}
